\section{Анализ предметной области}

\subsection{Иерархия памяти в современных вычислительных системах}

Производительность современных компьютеров определяется не только скоростью процессора, но и эффективностью работы подсистемы памяти. Процессоры работают на частотах, на порядки превышающих скорость доступа к основной оперативной памяти (DRAM). Чтобы сгладить эту разницу, используется многоуровневая иерархия кэш-памяти.

\begin{figure}[h]
\centering
\includegraphics[width=0.6\textwidth]{memory_hierarchy.png} % Предполагается наличие схемы
\caption{Типичная иерархия памяти}
\label{fig:mem}
\end{figure}

\begin{itemize}
    \item \textbf{Регистры процессора:} Самый быстрый тип памяти, расположенный непосредственно на ядре процессора. Время доступа — 1 такт. Объем крайне мал.
    \item \textbf{Кэш L1 (Level 1):} Небольшая (десятки КБ) статическая память (SRAM), индивидуальная для каждого ядра. Делится на кэш инструкций и кэш данных. Время доступа — несколько тактов. Хранит данные, которые используются в данный момент.
    \item \textbf{Кэш L2 (Level 2):} Более медленная и объемная (сотни КБ — мегабайты) SRAM-память. Может быть индивидуальной для ядра или общей для нескольких ядер. Время доступа — десятки тактов.
    \item \textbf{Кэш L3 (Level 3):} Самый большой и медленный кэш (десятки МБ), общий для всех ядер процессора. Служит последним рубежом перед обращением к ОЗУ. Время доступа — десятки-сотни тактов.
    \item \textbf{Оперативная память (DRAM):} Основная память системы (гигабайты). Значительно медленнее кэшей. Время доступа — сотни тактов.
    \item \textbf{Постоянные накопители (SSD/HDD):} Хранят данные постоянно. Время доступа на порядки медленнее ОЗУ (микросекунды для SSD, миллисекунды для HDD).
\end{itemize}

Эффективность работы этой иерархии основана на \textbf{принципе локальности}, который гласит, что программы имеют тенденцию повторно использовать данные и инструкции, к которым они недавно обращались (временная локальность), а также обращаться к данным, расположенным в памяти по соседству (пространственная локальность). Понимание задержек и пропускной способности каждого уровня иерархии является ключевым для написания производительного кода.

\subsection{Модели параллелизма и конкурентности}

С ростом числа ядер в процессорах параллельное программирование стало стандартом. Существует несколько моделей организации параллельных вычислений:

\begin{itemize}
    \item \textbf{Многопоточность с общей памятью (Shared-memory multithreading):} Несколько потоков выполнения работают в одном адресном пространстве и имеют доступ к общим данным. Это наиболее распространенная модель в современных ОС. Основная сложность — синхронизация доступа к общим данным для предотвращения состояний гонки. Для этого используются примитивы синхронизации: мьютексы, семафоры, атомарные операции.
    \item \textbf{Модель передачи сообщений (Message Passing):} Параллельные процессы имеют изолированную память и взаимодействуют исключительно путем отправки и получения сообщений. Ярким примером является стандарт MPI (Message Passing Interface), используемый в высокопроизводительных вычислениях (HPC). Эта модель позволяет избежать многих проблем с синхронизацией, но может вносить накладные расходы на сериализацию и пересылку данных.
    \item \textbf{Модель акторов:} Вариация модели передачи сообщений, где каждый «актор» — это легковесный процесс со своим состоянием, который может только получать сообщения и отправлять их другим акторам. Язык Go с его горутинами и каналами реализует схожую философию: «Не обменивайтесь данными через общую память; обменивайтесь памятью через передачу данных».
\end{itemize}

Тестирование производительности примитивов синхронизации и механизмов коммуникации (каналов) позволяет оценить, насколько эффективно можно распараллелить задачу на данной системе с использованием конкретного языка программирования.

\subsection{Обзор существующих средств анализа и профилирования}

На рынке существует множество инструментов для анализа производительности, каждый со своей нишей.

\begin{itemize}
    \item \textbf{Системные утилиты ОС:} Такие как \texttt{top}, \texttt{htop}, \texttt{iostat} в Linux/macOS и «Диспетчер задач»/«Монитор ресурсов» в Windows. Они предоставляют высокоуровневую картину загрузки системы в реальном времени, но не позволяют проводить глубокий синтетический анализ или измерять специфические параметры, такие как задержка кэша.
    \item \textbf{Профилировщики от производителей ЦП:} Intel VTune Profiler, AMD uProf. Это мощнейшие инструменты, использующие доступ к аппаратным счетчикам производительности (Performance Monitoring Units, PMU). Они могут отслеживать такие события, как промахи в кэше, неверно предсказанные ветвления и т.д. Их главный недостаток — сложность и привязка к конкретному производителю.
    \item \textbf{Синтетические бенчмарки:} Наборы тестов, такие как Geekbench, Phoronix Test Suite, AIDA64. Они популярны для сравнения общей производительности систем, но часто представляют собой «черный ящик», не раскрывая методологию своих тестов и не позволяя анализировать отдельные низкоуровневые компоненты.
    \item \textbf{Инструменты профилирования в языках программирования:} Например, встроенный в Go профилировщик pprof. Он отлично подходит для анализа производительности конкретного Go-приложения (поиск «горячих» функций, анализ использования памяти), но не предназначен для анализа аппаратной платформы в целом.
\end{itemize}

Анализ показывает, что существует незанятая ниша для инструмента, который бы сочетал простоту и наглядность системных утилит с возможностью проведения понятных и воспроизводимых синтетических тестов для ключевых аппаратных компонентов, будучи при этом кросс-платформенным и не привязанным к конкретному производителю.

